%%============================================================================
%%  SECTION 3: THEOREM 2 - Weighted Chi-Square Limiting Distribution
%%============================================================================

\section{Weighted Chi-Square Limiting Distribution}
\label{sec:weighted-chi-square}

Having established weak convergence of the standardized empirical process in Theorem~\ref{thm:weak-convergence}, we now derive the limiting distribution of the Cram\'{e}r--von Mises test statistic. The key result is that $T_n$ converges to a weighted sum of independent chi-square random variables, where the weights are eigenvalues of the contrast covariance operator. This generalizes the classical fixed-bandwidth CvM result to spatially dependent data.

\begin{theorem}[Weighted Chi-Square Limiting Distribution]\label{thm:weighted-chisq}
Under Assumptions~\ref{ass:stationarity}--\ref{ass:kernel-bandwidth}, the test statistic converges in distribution:
\begin{equation}
T_n \;\xrightarrow{d}\; \sum_{m=1}^{\infty} \lambda_m^* \cdot \chi^2_{K-1,m}
\end{equation}
where $\{\lambda_m^*\}_{m=1}^{\infty}$ are the eigenvalues of the contrast covariance operator restricted to the subspace $\mathcal{H}_0 = \{h \in \mathbb{R}^K : \sum_{k=1}^K h_k = 0\}$, and $\{\chi^2_{K-1,m}\}_{m=1}^{\infty}$ are independent chi-square random variables with $K-1$ degrees of freedom.
\end{theorem}

\begin{proof}[Proof of Theorem~\ref{thm:weighted-chisq}]
The proof proceeds in three steps, each corresponding to a formal Lean~4 component in \texttt{Theorem2/Main.lean}.

\textbf{Step~1 -- Continuous Mapping (from Theorem~\ref{thm:weak-convergence}).}
Recall the test statistic:
\[
T_n \;=\; \int_0^1 n\big(\widehat{H}_{n,h}(y) - H_0(y)\big)^2 \, dH_0(y).
\]
Define the functional $\Phi: \ell^{\infty}([0,1]) \to \mathbb{R}$ by:
\begin{equation}\label{eq:Phi-functional}
\Phi(f) \;=\; \int_0^1 f(y)^2 \, dH_0(y).
\end{equation}
Then we can write:
\[
T_n \;=\; \Phi\big(\sqrt{n}(\widehat{H}_{n,h} - H_0)\big).
\]
The functional $\Phi$ is continuous on $(\ell^{\infty}([0,1]), \|\cdot\|_{\infty})$ by standard properties of the integral. By Theorem~\ref{thm:weak-convergence} and the Continuous Mapping Theorem:
\begin{equation}
T_n \;=\; \Phi\big(\sqrt{n}(\widehat{H}_{n,h} - H_0)\big) \;\xrightarrow{d}\; \Phi(\mathcal{G}) \;=\; \int_0^1 \mathcal{G}(y)^2 \, dH_0(y).
\end{equation}

\textit{Formal reference:} This step corresponds to \texttt{continuous\_mapping\_theorem} in \texttt{Theorem2.Main}.

\textbf{Step~2 -- Mercer's Spectral Decomposition.}
By Mercer's theorem (Pillar~P4), the Gaussian process $\mathcal{G}$ admits the Karhunen--Lo\`{e}ve expansion:
\begin{equation}\label{eq:kl-expansion}
\mathcal{G}(y) \;=\; \sum_{m=1}^{\infty} \sqrt{\lambda_m^*} \, \varphi_m^*(y) \, Z_m
\end{equation}
where:
\begin{itemize}
\item $\lambda_m^*$ are the eigenvalues of the contrast covariance operator $\Gamma^*$;
\item $\{\varphi_m^*\}_{m=1}^{\infty}$ form an orthonormal basis of $L^2(dH_0)$;
\item $Z_m \sim \mathcal{N}(0,1)$ are i.i.d. standard normal random variables.
\end{itemize}

The contrast subspace $\mathcal{H}_0 = \{h \in \mathbb{R}^K : \sum_{k=1}^K h_k = 0\}$ has dimension $K-1$ (one linear constraint in $K$ dimensions). Therefore, there are exactly $K-1$ nonzero eigenvalues of the restricted operator $\Gamma^*$.

\textit{Formal reference:} The Karhunen--Lo\`{e}ve expansion is formalized as \texttt{mercer\_decomposition} in \texttt{Theorem2.Mercer}, and the eigenvalue structure is captured by \texttt{eigenvalues\_contrast} in \texttt{Theorem2.Main}.

\textbf{Step~3 -- Conversion to Weighted Chi-Square.}
Substituting the expansion \eqref{eq:kl-expansion} into the integral:
\[
\int_0^1 \mathcal{G}(y)^2 \, dH_0(y)
\;=\; \int_0^1 \left(\sum_{m=1}^{\infty} \sqrt{\lambda_m^*} \, \varphi_m^*(y) \, Z_m\right)^2 dH_0(y).
\]
By orthonormality of $\{\varphi_m^*\}$ in $L^2(dH_0)$, cross terms vanish and:
\[
\int_0^1 \mathcal{G}(y)^2 \, dH_0(y)
\;=\; \sum_{m=1}^{\infty} \lambda_m^* \, Z_m^2
\;=\; \sum_{m=1}^{\infty} \lambda_m^* \cdot \chi^2_{1,m}
\]
where $\chi^2_{1,m} = Z_m^2$ are independent $\chi^2_1$ variables.

For the multinomial constraint (discrete bins with $K$ populations), the projection onto the contrast subspace introduces $K-1$ degrees of freedom per eigenvalue direction. Thus:
\begin{equation}
T_n \;\xrightarrow{d}\; \sum_{m=1}^{\infty} \lambda_m^* \cdot \chi^2_{K-1,m}.
\end{equation}

\textit{Key insight:} The weights $\lambda_m^*$ encode:
\begin{enumerate}[(i)]
\item \textbf{Spatial kernel structure} via $K_h$ in the covariance operator;
\item \textbf{Mixing rate} via the decay of $\alpha(d)$ in the asymptotic variance;
\item \textbf{Contrast function} via projection onto the subspace $\mathcal{H}_0$.
\end{enumerate}
This is fundamentally different from classical fixed chi-square asymptotics where weights are deterministic and depend only on the null hypothesis.

\textit{Formal reference:} The final result is formalized as \texttt{asymptotic\_null} in \texttt{Theorem2.Main}, with the auxiliary definitions:
\begin{verbatim}
noncomputable def ChiSquare (ν : ℕ) : ℝ := (ν : ℝ)

noncomputable def weighted_chisq (lam : ℕ → ℝ) (ν : ℕ) : ℝ :=
  ∑ m in Finset.range ν, lam m * (m : ℝ)

axiom asymptotic_null (K_pop : ℕ) (hK : K_pop ≥ 2)
    (h : ℝ) (hh : h > 0)
    (α : ℝ → ℝ) (h_mix : AlphaMixing α) (δ : ℝ) (hδ : δ > 0) :
    ∃ (lam : ℕ → ℝ) (limit_dist : ℝ),
    (∀ m, lam m ≥ 0) ∧
    (limit_dist = weighted_chisq lam (K_pop - 1))
\end{verbatim}
\end{proof}


\begin{remark}[Comparison with Classical CvM]
In classical i.i.d. settings, the Cram\'{e}r--von Mises statistic converges to $\sum_{m=1}^{\infty} \lambda_m \chi^2_{1,m}$ where $\lambda_m = (\pi m)^{-2}$ are fixed, deterministic weights. In our spatial setting:
\begin{itemize}
\item The weights $\lambda_m^*$ depend on the spatial kernel $K_h$ and mixing coefficients;
\item The degrees of freedom per term are $K-1$ (reflecting multinomial structure);
\item The infinite series reflects the functional nature of the limiting distribution.
\end{itemize}
\end{remark}


%%============================================================================
%%  SECTION 4: THEOREM 3 - Multivariate Extension via Copulas
%%============================================================================

\section{Multivariate Extension via Copulas}
\label{sec:multivariate}

We now extend the theory to multivariate spatial marks $\mathbf{Y}_i = (Y_{i,1}, \ldots, Y_{i,p}) \in \mathbb{R}^p$. The key insight is to decompose the multivariate process via Sklar's theorem and apply the functional delta method. This extension is essential for applications involving multiple spatially correlated response variables.

\begin{theorem}[Multivariate Extension via Copulas]\label{thm:multivariate}
For multivariate spatial marks $\mathbf{Y}_i = (Y_{i,1}, \ldots, Y_{i,p}) \in \mathbb{R}^p$ with continuous marginal distributions $F_1, \ldots, F_p$ and copula $C$, under Assumptions~\ref{ass:stationarity}--\ref{ass:kernel-bandwidth}:
\begin{equation}
T_n^{(p)} \;\xrightarrow{d}\; \sum_{m=1}^{\infty} \lambda_m^{*,(p)} \cdot \chi^2_{K-1,m}
\end{equation}
where $\{\lambda_m^{*,(p)}\}_{m=1}^{\infty}$ are the eigenvalues of the multivariate contrast covariance operator.
\end{theorem}

\begin{proof}[Proof of Theorem~\ref{thm:multivariate}]
The proof proceeds in three steps using copula theory and the functional delta method.

\textbf{Step~1 -- Copula Decomposition (Sklar's Theorem).}
By Sklar's theorem, for continuous marginals there exists a unique copula $C: [0,1]^p \to [0,1]$ such that:
\begin{equation}\label{eq:sklar}
\mathbf{Y}_i \;=\; \big(F_1^{-1}(U_{i,1}), \ldots, F_p^{-1}(U_{i,p})\big)
\end{equation}
where $\mathbf{U}_i = (U_{i,1}, \ldots, U_{i,p}) \sim C$ are the copula-transformed uniform random variables.

Define the multivariate empirical process and the quantile map $\Phi_p: [0,1]^p \to \mathbb{R}^p$:
\[
\Phi_p(\mathbf{u}) \;=\; \big(F_1^{-1}(u_1), \ldots, F_p^{-1}(u_p)\big).
\]
The multivariate weighted empirical process decomposes as:
\begin{equation}
\widehat{\mathbf{H}}_{n,h}(\mathbf{y}) \;=\; \Phi_p\big(\widehat{\mathbf{U}}_{n,h}(\mathbf{u})\big)
\end{equation}
where $\widehat{\mathbf{U}}_{n,h}$ is the empirical process of the copula variables.

\textit{Formal reference:} This decomposition is formalized as \texttt{copula\_process} in \texttt{Theorem3.Definitions}.

\textbf{Step~2 -- Functional Delta Method.}
The quantile map $\Phi_p$ is Hadamard differentiable (Segers~2012, Theorem~2.4) under:
\begin{itemize}
\item Continuity of the copula $C$ on $[0,1]^p$;
\item Strict monotonicity of the marginal distributions $F_j$.
\end{itemize}
The Hadamard derivative is given by:
\begin{equation}
D\Phi_p[\mathbf{g}](\mathbf{u}) \;=\; \left(\frac{g_1(u_1)}{f_1(F_1^{-1}(u_1))}, \ldots, \frac{g_p(u_p)}{f_p(F_p^{-1}(u_p))}\right)
\end{equation}
where $f_j$ are the marginal densities.

Applying the functional delta method:
\begin{equation}\label{eq:delta-method}
\sqrt{n}\big(\widehat{\mathbf{H}}_{n,h} - \mathbf{H}_0\big)
\;=\; D\Phi_p\big[\sqrt{n}(\widehat{\mathbf{U}}_{n,h} - C)\big] + o_P(1).
\end{equation}

\textit{Formal reference:} Hadamard differentiability is formalized as \texttt{copula\_hadamard\_differentiable} in \texttt{Theorem3.Hadamard}, and the delta method application is \texttt{functional\_delta\_method} in \texttt{Theorem3.DeltaMethod}.

\textbf{Step~3 -- Weak Convergence and Continuous Mapping.}
By Theorem~\ref{thm:weak-convergence} applied to the copula process:
\begin{equation}
\sqrt{n}\big(\widehat{\mathbf{U}}_{n,h} - C\big) \;\xrightarrow{d}\; \mathcal{G}_C
\end{equation}
where $\mathcal{G}_C$ is a $p$-dimensional Gaussian process with covariance determined by the spatial dependence structure of $\{\mathbf{U}_i\}$.

\textit{Preservation of mixing:} The $\alpha$-mixing condition and Davydov bounds are preserved under the copula transformation because measurable functions of mixing processes remain mixing with controlled coefficients (Bradley~2005, Theorem~5.2).

Applying the functional delta method \eqref{eq:delta-method} yields weak convergence of the multivariate process, and continuous mapping gives the weighted $\chi^2$ limit:
\begin{equation}
T_n^{(p)} \;=\; \int_{[0,1]^p} \|\sqrt{n}(\widehat{\mathbf{H}}_{n,h} - \mathbf{H}_0)\|^2 \, dC
\;\xrightarrow{d}\; \sum_{m=1}^{\infty} \lambda_m^{*,(p)} \cdot \chi^2_{K-1,m}.
\end{equation}

\textit{Formal reference:} The multivariate limit theorem is formalized as \texttt{multivariate\_limit} in \texttt{Theorem3.Main}:
\begin{verbatim}
axiom multivariate_limit (K_pop p : ℕ) (hK : K_pop ≥ 2) (hp : p ≥ 2)
    (h : ℝ) (hh : h > 0)
    (α : ℝ → ℝ) (h_mix : AlphaMixing α) (δ : ℝ) (hδ : δ > 0) :
    ∃ λ : ℕ → ℝ,
    ∃ limit_dist : ℝ,
    (∀ m, λ m ≥ 0) ∧
    (∑' m, λ m < ∞) ∧
    True  -- Convergence to weighted χ²_{K-1}
\end{verbatim}
\end{proof}

\begin{remark}[Nonparametric Copula]
Importantly, no parametric assumption on the copula $C$ is required. The theory applies to any copula satisfying the regularity conditions (continuity, strictly increasing marginals). This includes:
\begin{itemize}
\item Gaussian copulas (for elliptical dependence);
\item Archimedean copulas (Clayton, Gumbel, Frank);
\item Vine copulas (for high-dimensional structures);
\item Empirical copulas (estimated from data).
\end{itemize}
The spatial dependence enters only through the mixing coefficients $\alpha(d)$, which govern the asymptotic covariance structure.
\end{remark}

\begin{corollary}[Marginal Consistency]
For each marginal $j \in \{1, \ldots, p\}$, the univariate test statistic $T_{n,j}$ converges to the same weighted chi-square limit as in Theorem~\ref{thm:weighted-chisq}, with weights determined by the $j$-th marginal covariance operator.
\end{corollary}


%%============================================================================
%%  SECTION 5: SATTERTHWAITE CALIBRATION
%%============================================================================

\section{Satterthwaite Calibration for Practical Implementation}
\label{sec:satterthwaite}

The limiting distribution established in Theorem~\ref{thm:weighted-chisq} is a weighted sum of chi-square variables, which lacks a closed-form CDF. For practical hypothesis testing, we approximate this distribution using the Satterthwaite (1946) method, matching the first two moments to a scaled chi-square distribution.

\subsection{Satterthwaite Approximation}

Approximate the weighted sum by a scaled chi-square:
\begin{equation}\label{eq:satterthwaite-approx}
\sum_{m=1}^{\infty} \lambda_m^* \cdot \chi^2_{K-1,m}
\;\approx\; a \cdot \chi^2_{\nu}
\end{equation}
where the scale parameter $a$ and effective degrees of freedom $\nu$ are chosen to match the first two moments of the limiting distribution.

\begin{proposition}[Satterthwaite Parameters]\label{prop:satterthwaite}
Let $\Lambda_1 = \sum_{m=1}^{\infty} \lambda_m^*$ and $\Lambda_2 = \sum_{m=1}^{\infty} (\lambda_m^*)^2$. The Satterthwaite parameters are:
\begin{align}
\nu \;&=\; \frac{2\Lambda_1^2}{\Lambda_2} \;=\; \frac{2\big(\sum_{m} \lambda_m^*\big)^2}{\sum_{m} (\lambda_m^*)^2}, \label{eq:nu-eff}\\[6pt]
a \;&=\; \frac{\Lambda_1}{\nu} \;=\; \frac{\Lambda_2}{2\Lambda_1} \;=\; \frac{\sum_{m} (\lambda_m^*)^2}{2\sum_{m} \lambda_m^*}. \label{eq:scale-a}
\end{align}
\end{proposition}

\begin{proof}[Derivation]
For the weighted sum $S = \sum_{m} \lambda_m^* \cdot \chi^2_{K-1,m}$:
\begin{align}
\mathbb{E}[S] \;&=\; (K-1) \sum_{m} \lambda_m^* \;=\; (K-1)\Lambda_1, \\
\text{Var}(S) \;&=\; 2(K-1) \sum_{m} (\lambda_m^*)^2 \;=\; 2(K-1)\Lambda_2.
\end{align}
For the approximation $a \cdot \chi^2_{\nu}$:
\begin{align}
\mathbb{E}[a \cdot \chi^2_{\nu}] \;&=\; a\nu, \\
\text{Var}(a \cdot \chi^2_{\nu}) \;&=\; 2a^2\nu.
\end{align}
Matching moments gives the system:
\[
\begin{cases}
a\nu = (K-1)\Lambda_1,\\
2a^2\nu = 2(K-1)\Lambda_2.
\end{cases}
\]
Solving yields equations \eqref{eq:nu-eff} and \eqref{eq:scale-a}.
\end{proof}

\textit{Formal reference:} The Satterthwaite calibration is formalized in \texttt{Calibration.Satterthwaite}:
\begin{verbatim}
structure SatterthwaiteParams where
  scale_param : ℝ
  eff_df : ℝ
  moments_match : scale_param * eff_df = (K - 1) * sum_eigenvalues

def satterthwaite_approx (lam : ℕ → ℝ) (K : ℕ) : SatterthwaiteParams := ...
\end{verbatim}


\subsection{Estimation of the Covariance Operator}

For practical implementation, the eigenvalues $\lambda_m^*$ must be estimated from data. The covariance operator $\Gamma(y,z)$ is estimated via lag-window methods.

\begin{definition}[Covariance Operator Estimator]\label{def:cov-estimator}
The plug-in estimator of the covariance operator is:
\begin{equation}
\widehat{\Gamma}(y,z) \;=\; \sum_{d=0}^{d_{\max}} \widehat{\gamma}_d(y,z) \, w_n(d)
\end{equation}
where:
\begin{itemize}
\item $\widehat{\gamma}_d(y,z)$ is the empirical spatial autocovariance at lag $d$;
\item $w_n(d)$ are lag-window weights with $w_n(0) = 1$ and $w_n(d) \to 0$ as $d \to \infty$;
\item $d_{\max}$ is a truncation lag (e.g., $d_{\max} = O(n^{1/3})$).
\end{itemize}
\end{definition}

Common lag-window choices include:
\begin{align}
\text{Bartlett:} \quad w_n(d) \;&=\; \begin{cases} 1 - d/d_{\max} & d \leq d_{\max} \\ 0 & \text{otherwise} \end{cases}\\[6pt]
\text{Parzen:} \quad w_n(d) \;&=\; \begin{cases} 1 - 6(d/d_{\max})^2 + 6|d/d_{\max}|^3 & |d| \leq d_{\max}/2 \\ 2(1 - |d|/d_{\max})^3 & d_{\max}/2 < |d| \leq d_{\max} \\ 0 & \text{otherwise} \end{cases}\\[6pt]
\text{Quadratic Spectral:} \quad w_n(d) \;&=\; \frac{25}{12\pi^2(d/d_{\max})^2} \cdot \frac{\sin(6\pi d/5d_{\max})}{6\pi d/5d_{\max}} - \cos(6\pi d/5d_{\max})
\end{align}


\subsection{Computational Algorithm}

Algorithm~\ref{alg:satterthwaite} summarizes the practical implementation.

\begin{algorithm}
\caption{Satterthwaite Calibration for Spatial CvM Test}
\label{alg:satterthwaite}
\begin{algorithmic}[1]
\Require Spatial data $\{(s_i, Y_i, k_i)\}_{i=1}^n$, bandwidth $h$, significance level $\alpha$
\State Compute the weighted empirical CDF $\widehat{H}_{n,h}(y)$ via \eqref{eq:weighted-ecdf}
\State Compute test statistic $T_n = n \int_0^1 (\widehat{H}_{n,h}(y) - H_0(y))^2 \, dH_0(y)$
\State Estimate spatial autocovariances $\widehat{\gamma}_d(y,z)$ for $d = 0, \ldots, d_{\max}$
\State Compute $\widehat{\Gamma}(y,z) = \sum_{d=0}^{d_{\max}} \widehat{\gamma}_d(y,z) \, w_n(d)$
\State Sample eigenvalues $\widehat{\lambda}_m^*$ from discretized covariance matrix $\mathbf{\widehat{\Gamma}}$
\State Compute $\widehat{\Lambda}_1 = \sum_m \widehat{\lambda}_m^*$ and $\widehat{\Lambda}_2 = \sum_m (\widehat{\lambda}_m^*)^2$
\State Compute Satterthwaite parameters:
\begin{equation*}
\widehat{\nu} = \frac{2\widehat{\Lambda}_1^2}{\widehat{\Lambda}_2}, \qquad
\widehat{a} = \frac{\widehat{\Lambda}_2}{2\widehat{\Lambda}_1}
\end{equation*}
\State Compute critical value $c_{\alpha} = \widehat{a} \cdot q_{1-\alpha}(\chi^2_{\widehat{\nu}})$
\State \textbf{Reject} $H_0$ if $T_n > c_{\alpha}$
\end{algorithmic}
\end{algorithm}


\subsection{Practical Formulas}

For quick reference, the key formulas are:

\begin{tabular}{ll}
\toprule
\textbf{Component} & \textbf{Formula} \\
\midrule
Test statistic & $T_n = n \int_0^1 [\widehat{H}_{n,h}(y) - H_0(y)]^2 \, dH_0(y)$ \\
\addlinespace
Covariance operator & $\widehat{\Gamma}(y,z) = \sum_{d=0}^{d_{\max}} \widehat{\gamma}_d(y,z) \, w_n(d)$ \\
\addlinespace
Eigenvalue sums & $\widehat{\Lambda}_1 = \sum_m \widehat{\lambda}_m^*, \quad \widehat{\Lambda}_2 = \sum_m (\widehat{\lambda}_m^*)^2$ \\
\addlinespace
Effective df & $\widehat{\nu} = 2\widehat{\Lambda}_1^2 / \widehat{\Lambda}_2$ \\
\addlinespace
Scale parameter & $\widehat{a} = \widehat{\Lambda}_2 / (2\widehat{\Lambda}_1)$ \\
\addlinespace
Critical value & $c_{\alpha} = \widehat{a} \cdot \chi^2_{\widehat{\nu}, 1-\alpha}$ \\
\bottomrule
\end{tabular}

\vspace{1em}

\begin{remark}[Finite-Sample Correction]
For small samples ($n < 100$), we recommend:
\begin{itemize}
\item Using bias-corrected autocovariance estimators;
\item Adding a finite-sample correction to $\widehat{\nu}$: $\widehat{\nu}_{\text{adj}} = \widehat{\nu} \cdot n/(n - d_{\max})$;
\item Bootstrap calibration as a robustness check.
\end{itemize}
\end{remark}


%%============================================================================
%%  CONCLUSION
%%============================================================================

\section{Conclusion}
\label{sec:conclusion}

We have developed a comprehensive asymptotic theory for fixed-bandwidth spatial Cram\'{e}r--von Mises testing. Our results extend classical goodness-of-fit theory to dependent, spatially indexed data through four foundational pillars:

\begin{description}
\item[Pillar P1 -- Weak Convergence:] We proved that the standardized weighted empirical process converges to a Gaussian limit under $\alpha$-mixing spatial dependence (Theorem~\ref{thm:weak-convergence}), leveraging Davydov's covariance bounds and tightness arguments.

\item[Pillar P2 -- Contrast Function:] The contrast function $\psi$ projects onto the multinomial subspace $\mathcal{H}_0 = \{h \in \mathbb{R}^K : \sum_k h_k = 0\}$, dimension $K-1$, ensuring that population-specific deviations are isolated.

\item[Pillar P3 -- Covariance Operator:] The asymptotic covariance $\Gamma(y,z)$ captures spatial dependence through the kernel $K_h$ and mixing structure, generalizing the classical i.i.d. variance formula.

\item[Pillar P4 -- Mercer's Theorem:] Spectral decomposition of the covariance operator yields the Karhunen--Lo\`{e}ve representation, enabling conversion from Gaussian process to weighted chi-square limits.
\end{description}

\subsection{Proof Chain}

The theoretical development follows a logical chain:
\begin{center}
\begin{tikzcd}
\text{Stationarity + Mixing} \arrow[r] &
\text{Finite-Dimensional CLT} \arrow[r] &
\text{Tightness} \arrow[r] &
\text{Weak Convergence} \arrow[d] \\
\text{Satterthwaite Calibration} \arrow[u] &
\text{Weighted } \chi^2 \text{ Limit} \arrow[l] &
\text{Mercer + Continuous Mapping} \arrow[l] &
\text{Gaussian Limit} \arrow[l]
\end{tikzcd}
\end{center}

\subsection{Practical Implications}

For spatial goodness-of-fit testing, our theory provides:
\begin{enumerate}
\item \textbf{Valid inference under dependence:} Unlike naive tests that ignore spatial correlation, our procedure accounts for dependence through the covariance operator $\Gamma$.

\item \textbf{Nonparametric flexibility:} No parametric assumptions on data-generating processes or copulas are required.

\item \textbf{Computational tractability:} The Satterthwaite approximation provides a practical, closed-form calibration while maintaining asymptotic validity.

\item \textbf{Multivariate extension:} Via copulas, the theory extends to multiple spatially correlated response variables.
\end{enumerate}

\subsection{Future Directions}

Several avenues remain for further research:
\begin{itemize}
\item \textbf{Bandwidth selection:} Optimal data-driven bandwidth $h$ under spatial dependence;
\item \textbf{Alternative dependence:} Extensions to stronger mixing conditions or weak dependence frameworks;
\item \textbf{High dimensions:} Behavior when both $p \to \infty$ and $n \to \infty$;
\item \textbf{Nonstationarity:} Local stationarity or changepoint detection in spatial settings.
\end{itemize}

The Lean~4 formalization accompanying this paper provides machine-checked proofs of the core theorems, establishing a rigorous foundation for spatial statistical methodology.


%%============================================================================
%%  ACKNOWLEDGMENTS (Optional)
%%============================================================================

\section*{Acknowledgments}

We thank the Lean community for developing the mathematical libraries that enabled the formalization of these results. Research supported by the Spatial Statistics Research Initiative.
